Table~\ref{tab:comparativo-trabalhos} summarizes the works of this chapter. None of the hardware works \textcolor{GREEN}{derives} a deadline for the inversion or the detection \textcolor{GREEN}{from} user mobility.
% FONTE: wang2021tlbdinsa, wang2023bdnsepa, gangarajaiah2017cholesky, li2019fdlama e kulkarni2021newton (sem km/h nem prazo; gangarajaiah L170 só fixa 20 M CD/s para banda larga; kulkarni L171 usa N_coh só para a banda de interconexão).

\begin{table}[htbp]
\caption[Comparison of the related works.]{Comparison of the related works. Metrics: R, resources or area; P, power; E, energy; T, throughput; D, delay or latency. --: no hardware implementation.}
\label{tab:comparativo-trabalhos}
\vspace{4pt} % a classe zera o espaço entre a caption e a tabela
\centering
\footnotesize
\renewcommand{\arraystretch}{1.15}
\setlength{\tabcolsep}{3pt}
\begin{tabular}{>{\raggedright\arraybackslash}p{2.9cm}>{\raggedright\arraybackslash}p{1.3cm}>{\raggedright\arraybackslash}p{2.0cm}>{\raggedright\arraybackslash}p{3.0cm}>{\raggedright\arraybackslash}p{2.8cm}>{\raggedright\arraybackslash}p{1.7cm}}
\toprule
\textbf{Work} & \textbf{Noise} & \textbf{Topology} & \textbf{Detection method} & \textbf{Platform} & \textbf{Metrics} \\
\midrule
\multicolumn{6}{l}{\textit{Decentralized algorithms}} \\
\cite{decentralized2017} & AWGN & Star & ADMM and CG & GPU cluster & T, D \\
\hline
\cite{jeon2019} & AWGN & Star (PD and FD) & ZF and LMMSE & Multiple GPUs & T, D \\
\hline
\cite{zheng2022} & AWGN & Chain & Recursive MMSE & -- & -- \\
\hline
\cite{chen2024grls} & AWGN & Daisy-chain & GRLS & -- & -- \\
\hline
\cite{rodriguezsanchez2020} & AWGN & Daisy-chain & Node-by-node update, no inversion & -- & -- \\
\hline
\cite{croisfelt2021} & AWGN & Daisy-chain & Kaczmarz, one equation per node, no inversion & -- & -- \\
\hline
\cite{li2023iotj} & AWGN & Star and daisy-chain & User-group detection, no inversion & -- & -- \\
\hline
\cite{zhuang2025} & AWGN & Star & LMMSE with imperfect CSI & -- & -- \\
\hline
\cite{zhao2021-globecom} & Colored & Daisy-chain & BCD-MMSE & -- & -- \\
\hline
\cite{zhao2025jsac} & Colored & Star and daisy-chain & BCD-MMSE, sDR-MMSE and cDR-MMSE & -- & -- \\
\midrule
\multicolumn{6}{l}{\textit{Inversion methods}} \\
\cite{li-tradeofs} & AWGN & Star & Explicit and implicit Cholesky & -- & -- \\
\hline
\cite{ingemarsson2016ldl} & -- & -- & LDL factorization, no square root & FPGA Virtex-7 & R, D \\
\hline
\cite{yu2015inversion} & -- & -- & SDF-SGR, no square root and no division & -- & -- \\
\midrule
\multicolumn{6}{l}{\textit{Hardware implementations}} \\
\cite{wang2023bdnsepa} & AWGN & Centralized & EP with block Neumann series & ASIC SMIC 65 nm & R, P, T \\
\hline
\cite{wang2021tlbdinsa} & AWGN & Centralized & Block partition, Cholesky inversion & FPGA Virtex-7 and ASIC TSMC 45 nm & R, P, T, D \\
\hline
\cite{gangarajaiah2017cholesky} & AWGN & Centralized & MF and ZF with Cholesky factorization & ASIC FD-SOI 28 nm & R, P, E, T \\
\hline
\cite{li2019fdlama} & AWGN & Star (FD) & LAMA, no inversion & FPGA Virtex-7 & R, T, D \\
\hline
\cite{kulkarni2021newton} & Not reported & Ring and star & Newton method & FPGA Virtex-7 & R, P, T, D \\
\midrule
This work & Colored & Daisy-chain & BCD-MMSE, Cholesky inversion of $\mathbf G_c$ & FPGA Alveo U50 (XCU50) & R, P, E, T, D \\
\bottomrule
\end{tabular}
\vspace{4pt}
\legend{Source: the author.}
\end{table}
% FONTE (linhas dos .md em referencias/):
%   Ruído: decentralized2017 L85; jeon2019 L75; li-tradeofs L45; zheng2022 L39; chen2024grls L31;
%     rodriguezsanchez2020 L63; croisfelt2021 L36 (sigma^2 I_M); li2023iotj L76; zhuang2025 L64;
%     zhao2021 e zhao2025 L8 (colored noise); wang2023bdnsepa L76; wang2021tlbdinsa L65; gangarajaiah2017cholesky L29;
%     li2019fdlama L34; kulkarni2021newton L52 ("channel noise", sem modelo).
%   Topologia: decentralized2017 L40 (central fusion node); jeon2019 L8 e L27 (PD e FD);
%     li-tradeofs L28 (feedforward); zheng2022 L8 (chain); chen2024grls L8; rodriguezsanchez2020
%     (título, daisy-chain); croisfelt2021 L20; li2023iotj L8 (star, extensão daisy-chain);
%     zhuang2025 (Fig. 1, clusters ligados a uma CU); zhao2025 L8 (star e daisy chain);
%     wang2023bdnsepa L610 (extensão descentralizada é trabalho futuro); gangarajaiah2017cholesky L8 (uma BS); kulkarni2021newton L40.
%   Plataforma e métricas: decentralized2017 L509-L511 (Cray XC30, Tabela III, latência e vazão);
%     jeon2019 L34 (multi-GPU, vazão e latência); wang2023bdnsepa L48 e resumo (SMIC 65 nm,
%     Gbps/W e Mbps/kGE); wang2021tlbdinsa e li2019fdlama: Tabela tab:hw_cholesky e texto deste
%     capítulo; gangarajaiah2017cholesky L17 e Tab. II (28 nm FD-SOI; gates, potência, energia, CD/s); kulkarni2021newton L217 e L278 (Virtex-7; latência, vazão e potência on-chip).
%     ingemarsson2016ldl L64 e Tab. II (latência); yu2015inversion L8 (sem implementação).
%     Este trabalho: METODOLOGIA/setup-hardware.tex e RESULTADOS.

None of the surveyed hardware works meets all three conditions of this work: (1)~colored noise
with a non-diagonal local covariance $\hat{\mathbf{R}}_{cc}$;
(2)~decentralized LMMSE equalization under the BCD-MMSE framework; and
(3)~FPGA implementation with measured resource metrics. The algorithm exists\textcolor{GREEN}{, but it has no hardware implementation}~\cite{zhao2025jsac}. This work fills this gap with a hardware analysis of the $M_c \times M_c$ Cholesky inversion of BCD-MMSE.
